\documentclass{beamer}
\usetheme{CambridgeUS}
\usepackage{amsmath}
\usepackage{amssymb}

\title{On Solving the Multiple Variable Gapped Longest Common Subsequence Problem}

\author{Marko Djukanović\inst{1,2, 6}  \and Nikola Balaban\inst{2}  \and
Christian Blum\inst{3}  \and 
Aleksandar Kartelj\inst{4}  \and
Sašo Džeroski\inst{5}  \and
Žiga Zebec\inst{6}
}

\institute{ $^1$University of Nova Gorica, Nova Gorica, Slovenia \\ 
$^2$Faculty of Natural Sciences and Mathematics, University of Banja Luka, Banja Luka, Bosnia and Herzegovina \\ 
$^3$Artificial Intelligence Research Institute (IIIA-CSIC), Barcelona, Spain \\ 
$^4$Faculty of Mathematics, University of Belgrade, Belgrade, Serbia \\
$^5$Jožef Stefan Institute, Ljubljana, Slovenia \ \\
$^6$Institute of Information Sciences (IZUM), Maribor, Slovenia \\ \vspace{0.3cm}

\footnotesize{-- \textcolor{blue}{Eurocast 2026: 20th International Conference on Computer Aided Systems Theory, February 23-27, 2026,  Las Palmas de Gran Canaria, Spain}--}
\centering
\includegraphics[width=90pt,height=50pt]{SMASH_CMYK-ENG-horizontal_on_white.pdf}~ \hfill
\includegraphics[width=80pt,height=50pt]{IFEEL\_1.png}~\hfill
\includegraphics[width=110pt,height=40pt]{CofundedbytheEU\_RGB\_NEG.png}
}

\begin{document}
\begin{frame}[plain]
    \maketitle
\end{frame}

\begin{frame}{Outline}
    
    \begin{itemize}
        \item Introduction \& Preliminaries \vspace{0.2cm}
        \item Problem Definition  \vspace{0.2cm}
        \item Graph state space  \vspace{0.2cm}
        \item Iterative multi-source Beam Search  \vspace{0.2cm}
        \item Experimental Evaluation
        \begin{itemize}
            \item General problem
            \item Special problem
        \end{itemize} \vspace{0.2cm}
        \item Conclusions
    \end{itemize}
\end{frame}

\begin{frame}{Introduction}
   \begin{itemize}
       \item Objects we deal with: sequences (strings) over finite alphabet 
       \begin{itemize}
       \item DNA/RNA over $\{\texttt{A}, \texttt{T}, \texttt{G}, \texttt{C/U}\}$
       \item Proteins over 20 (canonical) amino acids: $\{\texttt{A}, \texttt{C}, \texttt{D}, \texttt{E}, \texttt{P}, \texttt{Q} ...\}$ 
       \end{itemize} \vspace{0.3cm}
       \item Computational biology
       \begin{itemize}
          \item \textbf{One of central tasks:} sequence comparison, finding common motifs between sequences
          \item compare structurally but also semantically/functionality 
          \item sequence alignment problems 
       \end{itemize} \vspace{0.4cm}
       
       \item Subsequences: reveal structural similarities $\rightarrow$ \textcolor{blue}{\textbf{Longest commmon subsequence problem}} variants
   \end{itemize}
\end{frame}

\begin{frame}{Longest common subsequence problem (LCSP)}
    
    \begin{itemize}
      \item Basic problem in Computational biology
      \item Intensively solved over last 50 years
      \begin{itemize}
      \item theoretically as well as practically 
      \item Practically: many approximation algorothms, (meta-) heuristics, exact approaches, etc.
      \end{itemize}
    \end{itemize}
    
    \begin{definition}[LCSP]
        \textbf{Input}: Given a set of sequences $S=\{s_1, \ldots, s_m\}$ \\
        \textbf{Task}: Find a   subsequence $s$ which is \textcolor{blue}{common} for all sequences from $S$ of \textcolor{blue}{maximum} possible length. 
    \end{definition}
    \begin{example}
       Input: S=\{\texttt{\textcolor{red}{A}A\textcolor{red}{T}\textcolor{red}{T}G\textcolor{red}{C}}, \texttt{\textcolor{red}{A}\textcolor{red}{T}\textcolor{red}{T}A\textcolor{red}{C}}\} \\ 
       LCS solution: s=\textcolor{red}{ATTC}
    \end{example}
\end{frame}

\begin{frame}{Literature \& LCS Problem Variants }
Theoretically:
  \begin{itemize}
  \item When $m=2$ -- polynomially solvable (in $O(n^2)$): \textcolor{blue}{Dynamic programming}, Hunt-Schlimanski, ... 
  \item When $m$ arbitrary large -- $\mathcal{NP}$-hard:
   \begin{itemize}
   \item subject of interest within last 30 years: approximation approaches, meta-heuristics (ACO, \textcolor{blue}{Beam search}, ...), but also exact approaches (A$^*$, anytime approaches,...)
   \end{itemize}  
  \end{itemize}\vspace{0.3cm}
  
  Problem-related variants:
  \begin{itemize}
  		\item Arc-annotated LCS problem
 		 \item Constrained (restricted/imposed) LCS problem
  		\item Repetition-free, Longest filled LCS,...
  		\item \textbf{Gapped LCS problem}
  \end{itemize}
\end{frame}

\begin{frame}{The gaped LCS problem}
   \begin{definition}[A gap sequence]
       Given is a sequence $s$ and an assigned function 
       $G_s \colon \{1, \ldots, |s|\} \mapsto \mathbb{N}$. An ordered pair $(s, G_s)$ is called a \textbf{sequence with gaps}.
   \end{definition} \vspace{0.5cm}
   
   \begin{definition}[A gapped subsequence ]
         Sequence $\tilde{s}$ is a gapped subsequence of $(s, G_s)$ ifff 
         \begin{itemize}
         \item \textsc{s} is a subsequence of $s$
         \item the gapped constraint $G_s$ is fulfilled
         \end{itemize}
   \end{definition}
\end{frame}


\begin{frame}{Literature \& Motivation}
content...
\end{frame}


\end{document}
